\chapter{Integrative Methodology}
\label{ch:integrative-methodology}

\section{System Overview}

\systemname\ is designed as a modular pipeline rather than a monolithic model. The central assumption is that cross-media person identification benefits from decomposing the task into interpretable stages: detect the person, describe the person, represent the description, retrieve candidates, and then decide whether the person should be treated as a known or new identity. \Cref{fig:system-architecture} summarizes the current pipeline.

\begin{figure}[htbp]
  \centering
  \begin{tikzpicture}[
  node distance=0.9cm and 1.3cm,
  box/.style={draw, rounded corners, align=center, minimum width=3.2cm, minimum height=1cm, fill=blue!6},
  store/.style={draw, cylinder, shape border rotate=90, aspect=0.25, align=center, minimum height=1.2cm, minimum width=1.6cm, fill=green!8},
  arrow/.style={-{Latex[length=2.2mm]}, thick}
]
\node[box] (input) {Image / Video\\Frame};
\node[box, right=of input] (detector) {Module A\\Visual Extractor};
\node[box, right=of detector] (vlm) {Module B\\VLM Attribute Extractor};
\node[box, below=of vlm] (vectorizer) {Module C\\Dense + Sparse Vectorizer};
\node[box, left=of vectorizer] (matcher) {Module D\\Identity Matcher};
\node[store, below=of matcher] (db) {ChromaDB\\Vector Store};
\node[box, left=of matcher] (output) {Identity Decision\\and Ranked Matches};

\draw[arrow] (input) -- (detector);
\draw[arrow] (detector) -- (vlm);
\draw[arrow] (vlm) -- (vectorizer);
\draw[arrow] (vectorizer) -- (matcher);
\draw[arrow] (matcher) -- (output);
\draw[arrow] (vectorizer) |- (db);
\draw[arrow] (db) -- (matcher);
\end{tikzpicture}

  \caption{Current \systemname\ architecture. A person crop is detected, described by a VLM, vectorized into dense and sparse representations, and matched against a persistent vector database.}
  \label{fig:system-architecture}
\end{figure}

\section{Module A: Visual Extraction}

The first module uses a YOLO-family detector to localize person instances in an image or sampled video frame \citep{redmon2016you,jocher2023ultralytics}. Each detection is represented by a bounding box and confidence score. The prototype filters detections below a minimum bounding-box size and computes a simple quality score from detection confidence, relative crop size, crop sharpness, and centrality. For still-image processing, the highest-quality person crop is used as the downstream input. For video processing, detections are associated over time using intersection-over-union matching to produce lightweight trajectories.

The extractor deliberately avoids over-specialized assumptions about the source media. An uploaded still image, a cropped video frame, and a sampled surveillance frame are all reduced to the same intermediate object: a person crop with metadata about source path, bounding box, detection confidence, and quality score. This unification makes the rest of the pipeline media-agnostic.

\section{Module B: VLM Attribute Extraction}

The second module calls a VLM to convert a person crop into a structured JSON object. The prompt asks for attributes such as gender, age group, body build, topwear colour, topwear type, pattern description, bottomwear, shoes, bag, glasses, hat, mask, hair style, and posture. The response parser normalizes attribute keys, removes empty or unknown values, and attempts to repair malformed JSON when necessary. The current implementation supports cloud APIs and Aliyun DashScope-compatible Qwen-VL models, with environment-variable configuration for API keys.

The prompt is intentionally detailed for clothing patterns because preliminary tests showed that logos, stripes, and graphic prints can be more discriminative than broad colour categories. The VLM response is treated as uncertain evidence rather than ground truth. For this reason, the system stores the raw response, the cleaned attributes, and the source metadata for later inspection.

\section{Module C: Hybrid Vectorization}

The third module transforms the structured attributes into two complementary vector forms. A dense vector is generated by converting key--value attributes into a natural-language description and encoding that description with a sentence embedding model. A sparse vector is generated by mapping observed attribute keys to a persistent dynamic registry. The dense vector captures semantic similarity between descriptions, while the sparse vector preserves explicit overlap in attribute dimensions.

The current dense encoder uses a BGE-family embedding model, selected because it offers strong multilingual sentence representation and can handle Chinese attribute values in the prototype configuration \citep{xiao2023cpack}. The sparse registry is stored on disk so that attribute dimensions remain stable across sessions. This matters for incremental experiments: a database created today should remain interpretable after new attributes are observed tomorrow.

\section{Module D: Identity Matching}

The matcher retrieves candidate records from ChromaDB using dense vector similarity \citep{chroma2024}. It then recomputes a hybrid score:

\begin{equation}
S(q, c) = w_d S_d(q, c) + w_s S_s(q, c) + w_f S_f(q, c),
\end{equation}

\noindent where $S_d$ is dense cosine similarity, $S_s$ is sparse Jaccard-style similarity, and $S_f$ is reserved for a future face embedding component. In the current phase, $w_f=0$, and the default weights emphasize dense semantic similarity while retaining sparse attribute evidence. If the best candidate score is above a threshold, the query is assigned to the existing identity; otherwise, a new identity identifier is created.

The threshold is configured rather than learned in the current prototype. This is appropriate for the present phase because the evaluation dataset is still being constructed. Once a larger dataset is available, the threshold can be selected through validation curves and reported with precision--recall trade-offs.

\section{Dataset Creation Workflow}

The most recent system extension adds a GUI-supported dataset creation workflow. After video upload and trajectory extraction, the user can select one or more person tracks, choose sampling and crop parameters, and export a reusable dataset. The export includes cropped images, a manifest, positive same-identity pairs, negative cross-identity pairs, and a compressed ZIP archive. \Cref{fig:dataset-pipeline} shows this process.

\begin{figure}[htbp]
  \centering
  \begin{tikzpicture}[
  node distance=0.8cm and 1cm,
  box/.style={draw, rounded corners, align=center, minimum width=2.8cm, minimum height=0.9cm, fill=orange!7},
  file/.style={draw, align=center, minimum width=2.6cm, minimum height=0.8cm, fill=gray!8},
  arrow/.style={-{Latex[length=2.2mm]}, thick}
]
\node[box] (video) {Uploaded\\Video};
\node[box, right=of video] (tracks) {YOLO-based\\Track Extraction};
\node[box, right=of tracks] (select) {Manual Track\\Selection};
\node[box, below=of select] (sample) {Frame Sampling\\and Crop Export};
\node[file, left=of sample] (manifest) {manifest.csv\\manifest.jsonl};
\node[file, below=of sample] (pairs) {pairs.csv\\positive/negative};
\node[file, right=of sample] (zip) {Dataset ZIP\\for Overleaf / tests};

\draw[arrow] (video) -- (tracks);
\draw[arrow] (tracks) -- (select);
\draw[arrow] (select) -- (sample);
\draw[arrow] (sample) -- (manifest);
\draw[arrow] (sample) -- (pairs);
\draw[arrow] (sample) -- (zip);
\end{tikzpicture}

  \caption{GUI-supported video-to-testset workflow. The exported dataset is designed for later evaluation scripts and manual inspection.}
  \label{fig:dataset-pipeline}
\end{figure}

\section{Software Structure}

The codebase has been reorganized around a Python package and clear project-level directories. The core package contains configuration loading, application orchestration, CLI logic, algorithm modules, storage wrappers, and dataset export helpers. Applications such as the Streamlit GUI live in \texttt{apps/}; research scripts live in \texttt{experiments/}; generated datasets and runtime databases are excluded from version control by default.

\begin{table}[htbp]
  \centering
  \caption{Current repository organization relevant to the dissertation artifact.}
  \label{tab:repo-structure}
  \begin{tabularx}{\textwidth}{lX}
    \toprule
    Directory & Responsibility \\
    \midrule
    \texttt{crossmedia\_pid/} & Importable package: configuration, pipeline orchestration, core modules, storage, dataset export. \\
    \texttt{apps/} & Streamlit GUI entry points. \\
    \texttt{experiments/} & Evaluation and exploratory scripts. \\
    \texttt{experiments/generated\_datasets/} & GUI-exported crop datasets, ignored by version control. \\
    \texttt{configs/} & Local configuration and shareable example configuration. \\
    \texttt{docs/dissertation/} & Overleaf-compatible multi-paper dissertation draft. \\
    \bottomrule
  \end{tabularx}
\end{table}

\section{Evaluation Metrics}

The final dissertation evaluation will use pair-level and retrieval-level metrics. Pair-level evaluation treats each crop pair as a binary same-person or different-person decision. For a threshold $\tau$, a pair is predicted positive if its score is at least $\tau$. Precision, recall, and F1 are defined in the standard way:

\begin{equation}
\mathrm{Precision} = \frac{TP}{TP + FP}, \quad
\mathrm{Recall} = \frac{TP}{TP + FN}, \quad
\mathrm{F1} = \frac{2PR}{P + R}.
\end{equation}

Retrieval-level evaluation will use top-$k$ accuracy and mean average precision where a gallery is available. The GUI-generated \texttt{pairs.csv} file supports pair-level analysis immediately, while the \texttt{manifest.csv} file supports later conversion into query/gallery retrieval splits.

\section{Experimental Design}

The dissertation will distinguish three experiment types. First, \textbf{component ablations} compare dense-only, sparse-only, and hybrid scores. Second, \textbf{workflow experiments} evaluate whether GUI-generated datasets produce consistent labels and useful failure analysis metadata. Third, \textbf{cross-media extension experiments} test whether text descriptions can retrieve the correct person crops. This staged design avoids relying on a single headline score and instead evaluates the system as both a matching method and a research artifact.
